\documentclass[10pt,a4paper]{amsart}
\usepackage[margin=19mm]{geometry}
\usepackage{amsmath,amssymb,mathtools}
\usepackage[scr=rsfs,cal=euler]{mathalfa}
\usepackage{xfrac}
\usepackage[unicode,hidelinks,bookmarks=false]{hyperref}
\newcommand{\ie}{i.e.\ }
\newcommand{\cf}{cf.\ }
\newcommand{\resp}{respectively\ }
\newcommand{\Subsection}{\subsection}
\providecommand{\nobreakdash}{\nobreak\hbox{-}}
\setlength{\emergencystretch}{2em}
\multlinegap0pt
\usepackage{xcolor}
\usepackage{enumerate}
\usepackage{amssymb}
\usepackage{bm}
\newtheorem{proposition}{Proposition}
\newcommand{\G}{\mathrm{G}}
\newcommand{\h}{H}%{\textcolor{red}{H}}
\newcommand{\qthreep}{\mathbb{Q}^3_+}
\title{Bernstein-type Theorems \\ for constant mean curvature surfaces \\ in the three-dimensional light cone}
\author{Shintaro Akamine, Wonjoo Lee, Seong-Deog Yang}
\date{Source: arXiv:2603.17727v1 (CC BY 4.0)}
\begin{document}
\maketitle

\noindent\emph{Source and licence.} Bernstein-type Theorems \\ for constant mean curvature surfaces \\ in the three-dimensional light cone. Version arXiv:2603.17727v1 (CC BY 4.0), distributed by arXiv under the Creative Commons Attribution 4.0 International licence, http://creativecommons.org/licenses/by/4.0/, as recorded in the arXiv licence metadata for this exact version and shown on the abstract page https://arxiv.org/abs/2603.17727v1. Full licence notice: \texttt{licenses/arxiv-2603.17727-CC-BY-4.0.txt}.\par\smallskip

\noindent\textit{Editorial scope.} Counted unit: the article's proposition Prop:202601150810PM (source0.tex lines 436--439). The article's complete original proof of it is delivered with it (source0.tex lines 441--473, one \texttt{proof} environment of 33 lines). No mathematics is written, repaired or re-derived: the statement and the proof are the article's own lines, in the article's own notation, and no retained line is substituted or re-flowed.  Retained uncounted as the source-local context the proof needs: lines 363--365; lines 386--392.  Nothing was omitted from any retained span, so the package carries no omission record.  No retained line contains a citation, so the wrapper carries no bibliography and no bibliography entry.  No retained line contains a figure environment, an image-inclusion command or drawing code, so no image is delivered and the package declares no asset dependency.  Editorial formatting only, in addition: an AMS \texttt{amsart} preamble that restates the article's own declarations and macros used by the retained lines, namely the environments \texttt{proposition} on separate counters, together with the standard packages those lines need.  The retained environments print in this document as Proposition 1 in order of appearance, whereas the article prints the same items with the numbering of its own full document.  The complete native source of the article as distributed in the arXiv e-print for this exact version is bundled unchanged as \texttt{sources/196644/source0.tex}.
\begin{align}\label{Eq:202507142137}
	\langle \G, \G \rangle = \langle \G, X_u \rangle = \langle \G, X_v \rangle = 0,\qquad \langle \G, X \rangle = 1,
\end{align}

\begin{equation}\label{Eq:202601281208PM}
	\begin{cases}
		X_{zz} = Q X + 2 \omega_z X_z,                                      \\
		X_{z\bar{z}} = \frac{1}{2}He^{2\omega}X - \frac{1}{2}e^{2\omega}\G, \\
		\G_z = -HX_z - 2Qe^{-2 \omega}X_{\bar{z}}
	\end{cases}
\end{equation}

\begin{proposition}\label{Prop:202601150810PM}
	Any connected totally umbilic surface in $\qthreep$ is a (part of a) sphere, horosphere, or a hypersphere.
	In particular, any connected totally geodesic surface in $\qthreep$ is a (part of a) horosphere.
\end{proposition}\label{Prop:202601150642AM}

\begin{proof}
	Suppose $X$ is totally umbilic. Then, it follows from \eqref{Eq:202601281208PM} that for any $(u,v)$
	$$
		G_u(u,v) = -H(u,v) X_u(u,v), \qquad
		G_v(u,v) = -H(u,v) X_v(u,v)
	$$
	for the mean curvature function $H(u,v) $.
	From $G_{uv} = G_{vu}$, we see that $H_u=H_v=0$ for all $(u,v)$, so $H$ is constant, say $\h$.

	Suppose that $\h = 0$. Then $G$ is constant, say, $G(u,v)=G_0$. Then, from \eqref{Eq:202507142137}, we see that $\langle X, G_0 \rangle = 1$,
	hence $X$ lies in a lightlike hyperplane.

	Suppose that $\h \not= 0$. Then
	$
		G(u,v) = -\h ( X(u,v) - \vec{d} )
	$
	for some constant vector $\vec{d}$.
	Hence
	$$
		0 = \langle G, G \rangle = \h^2 \langle X-\vec{d}, X-\vec{d} \rangle = \h^2 ( \langle X, X \rangle - 2 \langle X, \vec{d} \rangle + \langle \vec{d},\vec{d} \rangle).
	$$
	Therefore,
	$
		\langle X, \vec{d} \rangle = \tfrac{1}{2} \langle \vec{d}, \vec{d} \rangle.
	$
	So $X$ lies in a hyperplane.
	Note that $\vec{d} = X + \tfrac{1}{\h} G$ hence
	$
		\langle \vec{d}, \vec{d} \rangle = 2/\h.
	$
	Therefore, if $\h<0$, then $\vec{d}$ is timelike and the hyperplane is spacelike.
	If $\h>0$, then $\vec{d}$ is spacelike and the hyperplane is timelike. Now the conclusion follows.
\end{proof}

\end{document}
